\section{Tối ưu bằng Pruning}

\subsection{Lựa chọn phương pháp pruning}

Đồ án ưu tiên structured channel pruning vì TensorRT có thể khai thác trực tiếp
kích thước tensor đã giảm. Điểm quan trọng của channel được tính bằng
\cbs{chuẩn trọng số/gradient/tiêu chí thực tế}. Các tầng đầu vào, detection head
và những nhánh có ràng buộc shape được bảo vệ hoặc cắt theo nhóm tương thích.

Để đánh giá lựa chọn này, một cấu hình unstructured pruning có thể được dùng làm
đối chứng. Nếu sparsity tăng nhưng latency không giảm, kết quả cho thấy sự khác
biệt giữa nén số tham số và tăng tốc thực tế.

\subsection{Xác định pruning ratio}

Ba mức khởi đầu là 20\%, 30\% và 40\%. Với mỗi mức, pipeline tạo một cấu hình
độc lập, ghi số channel còn lại, số tham số, FLOPs và checkpoint trước fine-tuning.
Nếu mô hình mất ổn định ở một mức, quá trình dừng sớm và ghi rõ tầng gây lỗi thay
vì tiếp tục tạo kết quả không hợp lệ.

\subsection{Fine-tuning}

Mô hình sau pruning được fine-tune từ \cbs{số epoch} epoch với learning rate
\cbs{giá trị}. Mask hoặc cấu trúc đã cắt được khóa để phần tử bị loại bỏ không
xuất hiện trở lại. Validation mAP được theo dõi theo epoch; checkpoint tốt nhất
được dùng cho đánh giá test và làm đầu vào cho QAT.

\subsection{Đánh giá mô hình sau pruning}

\begin{table}[H]
\centering
\caption{Kết quả các cấu hình pruning}
\label{tab:pruning-results}
\begin{tabularx}{\textwidth}{L{3cm}C{2cm}C{2.2cm}C{2.2cm}C{2cm}Y}
\toprule
\textbf{Cấu hình} & \textbf{Ratio} & \textbf{Params} & \textbf{\shortstack{mAP@50:\\95}} &
\textbf{Size} & \textbf{Nhận xét}\\
\midrule
Baseline & 0\% & \cbs{} & \cbs{} & \cbs{} & Đường chuẩn\\
Pruned-20 & 20\% & \cbs{} & \cbs{} & \cbs{} & \cbs{}\\
Pruned-30 & 30\% & \cbs{} & \cbs{} & \cbs{} & \cbs{}\\
Pruned-40 & 40\% & \cbs{} & \cbs{} & \cbs{} & \cbs{}\\
\bottomrule
\end{tabularx}
\end{table}

Đánh giá pruning không kết luận chỉ từ bảng trên mà tiếp tục đo TensorRT latency.
Cấu hình có ít tham số hơn nhưng tạo shape không thuận lợi cho kernel có thể
không đem lại speedup tương ứng.